% from ega4-1-fr.tex
\oldpage[IV-1]{230}

\begin{env}[1.3.10]
\label{IV.1.3.10-fr}
Si $B$ est une $A$-algèbre \emph{locale} essentiellement de type fini, $B$ est
de la forme $C_{\mathfrak q}$, où $C$ est une $A$-algèbre de type fini et
$\mathfrak q$ un idéal premier de $C$ (Bourbaki, \emph{Alg. comm.}, chap. II,
\S~5, prop. 11). Soit $\mathfrak p$ l'image réciproque dans $A$ de l'idéal
$\mathfrak q$~; si l'on pose $S=A-\mathfrak p$, $C_{\mathfrak q}$ est aussi un
anneau local en un idéal premier de $S^{-1}C$~; comme $S^{-1}C$ est une
algèbre de type fini sur $A_{\mathfrak p}=S^{-1}A$, on voit que $B$ est aussi
une $A_{\mathfrak p}$-algèbre \emph{essentiellement de type fini}, et en outre
l'homomorphisme $A_{\mathfrak p}\to B$ est \emph{local}.
\end{env}

\begin{proposition}[1.3.11]
\label{IV.1.3.11-fr}
Si $B$ est une $A$-algèbre locale essentiellement de type fini, $B$ est
$A$-isomorphe à une $A$-algèbre quotient d'une $A$-algèbre de la forme
$C_{\mathfrak q}$, où $C$ est une algèbre de polynômes
$A[T_1,\ldots,T_n]$ et $\mathfrak q$ un idéal premier de $C$.
\end{proposition}

\begin{proof}
En effet, par définition, $B$ est isomorphe à $C'_{\mathfrak q'}$, où $C'$ est
une $A$-algèbre de type fini et $\mathfrak q'$ un idéal premier de $C'$~; mais
$C'=C/\mathfrak b$, où $C=A[T_1,\ldots,T_n]$ et $\mathfrak b$ est un idéal de
$C$~; donc $\mathfrak q'=\mathfrak q/\mathfrak b$, où $\mathfrak q$ est un
idéal premier de $C$, et $C'_{\mathfrak q'}$ est isomorphe à
$C_{\mathfrak q}/\mathfrak bC_{\mathfrak q}$.
\end{proof}

\subsection{Morphismes localement de présentation finie}
\label{subsection:IV.1.4-fr}

\begin{env}[1.4.1]
\label{IV.1.4.1-fr}
Étant donné un anneau $A$, nous dirons qu'une $A$-algèbre $B$ (commutative)
est \emph{de présentation finie} si elle est isomorphe au quotient
$A[T_1,\ldots,T_n]/\mathfrak a$ d'une algèbre de polynômes sur $A$ par un
idéal \emph{de type fini} $\mathfrak a$ de $A[T_1,\ldots,T_n]$. Pour toute
$A$-algèbre (commutative) $A'$, $B\otimes_AA'$ est alors une $A'$-algèbre de
présentation finie, étant isomorphe au quotient de $A'[T_1,\ldots,T_n]$ par
l'image canonique dans cet anneau de $\mathfrak a\otimes_AA'$, qui est
évidemment un $A'[T_1,\ldots,T_n]$-module de type fini. Si $B$ est une
$A$-algèbre de présentation finie et $C$ une $B$-algèbre de présentation
finie, alors $C$ est une $A$-algèbre de présentation finie. En effet, soient
$B=A[S_1,\ldots,S_m]/\mathfrak a$ et
$C=B[T_1,\ldots,T_n]/\mathfrak b$, où $\mathfrak a$ (resp. $\mathfrak b$) est
un idéal de type fini de $A[S_1,\ldots,S_m]$ (resp.
$B[T_1,\ldots,T_n]$)~; l'anneau $B[T_1,\ldots,T_n]$ est isomorphe à
$A[S_1,\ldots,S_m,T_1,\ldots,T_n]/\mathfrak a'$, où $\mathfrak a'$ est
l'image canonique de
$\mathfrak a\otimes_AA[T_1,\ldots,T_n]$, et est donc un idéal de type fini de
$A[S_1,\ldots,S_m,T_1,\ldots,T_n]$~; l'idéal $\mathfrak b$ de
$B[T_1,\ldots,T_n]$ est de la forme $\mathfrak b'/\mathfrak a'$, où
$\mathfrak b'$ est un idéal de $A[S_1,\ldots,S_m,T_1,\ldots,T_n]$~; comme
$\mathfrak a'$ et $\mathfrak b'/\mathfrak a'$ sont des modules de type fini
sur $A[S_1,\ldots,S_m,T_1,\ldots,T_n]$, il en est de même de $\mathfrak b'$,
et $C$, isomorphe à
$A[S_1,\ldots,S_m,T_1,\ldots,T_n]/\mathfrak b'$, est donc bien une
$A$-algèbre de présentation finie. Notons enfin que si $A$ est
\emph{noethérien}, il en est de même de $A[T_1,\ldots,T_n]$, donc toute
$A$-algèbre de type fini est alors de présentation finie.
\end{env}

\begin{definition}[1.4.2]
\label{IV.1.4.2-fr}
Soient $f:X\to Y$ un morphisme de préschémas, $x$ un point de $X$, $y=f(x)$.
On dit que $f$ est \emph{de présentation finie en $x$} s'il existe un voisinage
ouvert affine $V$ de $y$ et un voisinage ouvert affine $U$ de $x$ tels que
$f(U)\subset V$ et que l'anneau $A(U)$ soit une $A(V)$-algèbre de présentation
finie. On dit que $f$ est \emph{localement de présentation finie} s'il est de
présentation finie en tout point de $X$.
\end{definition}

Si $Y$ est localement noethérien, il revient au même de dire que $f$ est
localement de type fini ou localement de présentation finie.

\begin{proposition}[1.4.3]
\label{IV.1.4.3-fr}
\begin{enumerate}
  \item[(i)] Tout isomorphisme local est localement de présentation finie.
  \item[(ii)] Si deux morphismes $f:X\to Y$, $g:Y\to Z$ sont localement de
  présentation finie, il en est de même de $g\circ f$.

